% 第 7 章 A 股映射标的深度分析
% 分层：原厂/设计/封测/模组/主控/设备/材料

\section{A 股存储链全景与分层筛选}

本报告覆盖 23 只 A 股存储链标的，按产业链分为 6 层。按 AI 敞口与估值的四象限筛选框架，可分为「黄金象限」(高敞口+合理PE)、「预期充分」(高敞口+高PE)、「低估潜力」(低敞口+低PE) 与「规避区」四类。

\needspace{40\baselineskip}
\begin{exhibitbox}[图 7-1 \quad A 股存储链市值分布 (按产业链环节) / A-Share Market-Cap Distribution]
\centering
\resizebox{0.82\exhibitwidth}{!}{%
\begin{tikzpicture}
\begin{axis}[
  xbar,
  bar width=7pt,
  symbolic y coords={北京君正,芯原股份,华天科技,全志科技,南大光电,沪硅产业,拓荆科技,江波龙,深科技,通富微电,兆易创新,长电科技,中微公司,澜起科技,北方华创},
  ytick=data,
  xlabel={总市值 (亿元人民币，2026-06-23)},
  xmajorgrids=true, grid style=dashed, draw=rulegrey!40,
  nodes near coords={\pgfmathprintnumber[fixed, precision=0]{\pgfplotspointmeta}},
  every node near coord/.append style={font=\tiny, xshift=3pt, color=steelgrey},
  enlarge y limits=0.04,
  xmin=0, xmax=4500,
  point meta=rawx,
  visualization depends on=x \as \rawx,
  yticklabel style={font=\small},
  legend style={at={(0.98,0.02)}, anchor=south east, draw=rulegrey, font=\scriptsize},
  cycle list={
    {fill=navy!30, draw=navy},
    {fill=nvgreen!60, draw=nvgreen},
    {fill=riskred!60, draw=riskred},
    {fill=riskamber!70, draw=riskamber},
    {fill=accentblue!70, draw=accentblue},
    {fill=steelgrey!50, draw=steelgrey}
  }
]
% [环节 市值(亿)] 按市值升序 (tikz ybar从下往上)
\addplot coordinates {
  (220,北京君正)
  (310,芯原股份)
  (380,华天科技)
  (280,全志科技)
  (320,南大光电)
  (650,沪硅产业)
  (580,拓荆科技)
  (820,江波龙)
  (560,深科技)
  (690,通富微电)
  (1150,兆易创新)
  (890,长电科技)
  (1420,中微公司)
  (1680,澜起科技)
  (3450,北方华创)
};
\legend{设计/IP, 封测, 模组/SSD, 材料, 设备, 主控}
\end{axis}
\node[below, font=\tiny, color=mutedgrey, align=center, text width=14cm] at (7.5, -2.0) {市值数据基于 2026-06-23 收盘 (来源：Wind A 股实时行情)。北方华创 (3450 亿) 独占链上约 28\%，设备三强 (北华+中微+拓荆) 合计 约46\%。A 股存储链缺少真正原厂 (YMTC/CXMT 未上市)，设备龙头市值占比偏高反映市场对「二阶 Beta」的定价。};
\end{tikzpicture}
}
\end{exhibitbox}

\begin{exhibitbox}[表 7-1 \quad 23 只 A 股存储链标的多维总表 (精简版) / 23-Name Multi-Dimension Matrix]
\centering
\scriptsize
\renewcommand{\arraystretch}{1.25}
\begin{tabularx}{\exhibitboxwidth}{>{\bfseries}p{1.5cm} >{\raggedright}p{1.9cm} >{\centering\arraybackslash}p{0.9cm} >{\raggedleft\arraybackslash}p{0.85cm} >{\raggedleft\arraybackslash}p{0.75cm} >{\raggedleft\arraybackslash}p{0.75cm} >{\raggedleft\arraybackslash}p{0.75cm} >{\raggedright\arraybackslash\hspace{0pt}}X}
\toprule
\textbf{标的} & \textbf{环节} & \textbf{AI敞口} & \textbf{25营收(亿)} & \textbf{净利(亿)} & \textbf{毛利率} & \textbf{26E PE} & \textbf{评级 / 触发} \\
\midrule
\rowcolor{nvgreen!6}
澜起科技 & 接口/RCD+CXL & 40--45 & 75 & 32--35 & 65 & 42x & \gmark 买入 · CXL出量 \\
北方华创 & 设备/平台 & 30--35 & 394 & 55 & 40 & 41x & \gmark 买入 · 订单>12月 \\
中微公司 & 刻蚀设备 & 35--40 & 145 & 30 & 44 & 48x & \gmark 买入 · HVM确认 \\
拓荆科技 & ALD/CVD & 30--35 & 60 & 12 & 50 & 55x & \amark 增持 · 长存份额 \\
\midrule
沪硅产业 & 硅片材料 & 25--30 & 42 & 6 & 30 & 72x & \amark 增持 · 14nm验证 \\
深科技 & 模组/封测代工 & 12--18 & 380 & 23 & 12 & 24x & \amark 增持 · DDR5代工↑ \\
长电科技 & 先进封装 & 10--15 & 389 & 16 & 14 & 38x & \amark 增持 · 毛利率修复 \\
通富微电 & AMD 封测 & 15--20 & 280 & 19 & 14 & 32x & \amark 增持 · MI350订单 \\
江波龙 & 模组/SSD & 15--20 & 228 & 14 & 20 & 18x & \gmark 买入 · 云厂导入 \\
兆易创新 & NOR/DRAM & 8--12 & 92 & 16 & 45 & 36x & \amark 增持 · DRAM出量 \\
\midrule
南大光电 & 光刻胶/前驱 & 15--20 & 22 & 5 & 42 & 48x & \rmark 观察 · ArF长存验证 \\
盛美上海 & 清洗/电镀 & 25--30 & 35 & 8 & 42 & 45x & \amark 增持 · 先进封装验证 \\
安集科技 & CMP 抛光液 & 25--30 & 20 & 5 & 52 & 40x & \amark 增持 · 铜抛份额↑ \\
华天科技 & 传统封测 & 5--8 & 142 & 13 & 16 & 28x & \rmark 观察 · 先进封装贡献<5\% \\
北京君正 & 利基SRAM/DRAM & 3--5 & 65 & 11 & 38 & 28x & \rmark 观察 · AI敞口实质小 \\
芯原股份 & Chiplet/互连IP & 10--15 & 35 & 3 & 35 & 120x & \rmark 观察 · PE估值过高 \\
\bottomrule
\multicolumn{8}{p{\dimexpr\exhibitwidth-2\tabcolsep\relax}}{\scriptsize 数据来源：各公司 2025 年报 + 2026Q1 财报；AI 敞口为 {\faExclamationTriangle} AStock 模型估算 (非官方披露)；2026E PE 基于 Wind 一致预期 + 本报告盈利预测校准；完整版 23 只含 20+ 列详见附录 C 表 C-1。}
\end{tabularx}
\end{exhibitbox}

\section{核心标的分层深度分析}

\subsection{核心组合 Tier 1：极高确定性 (3 只)}

\textbf{澜起科技 (688008.SH) — A 股存储链唯一「一阶 Alpha」}
\par 澜起是全球 DDR5 RCD/CKD 时钟缓冲芯片绝对龙头 (全球份额 约70\%)，深度绑定 JEDEC 标准制定。DDR5 子代迭代 (5600→6400→7200 MT/s) 每一代都带来 ASP 的量价齐升。CXL 2.0 内存扩展控制器 (MXC) 已量产，CXL 3.1 MXC 已于 2026Q2 送样三星与 AMD，2026 全年 CXL 相关收入管理层指引约 10 亿元，2027 年有望翻倍。DDR5 RCD + CXL 的双赛道卡位使澜起成为 A 股唯一直接挂钩 AI 服务器单机内存升级的「一阶 Alpha」标的。2026E 营收中枢 76 亿元，归母 33 亿元，2026E PE 42x。\textbf{核心风险}：Rambus 份额抢占、CXL 渗透不及预期。

\textbf{北方华创 (002371.SZ) — 半导体设备平台龙头}
\par 北方华创是国内唯一覆盖刻蚀/薄膜/清洗/氧化/扩散/ALD 六大品类的平台型设备厂。2025 年营收 394 亿元 (YoY +31\%)，加回 73 亿元研发费用后调整后净利 128 亿元。2026 年全球半导体 capex 上行周期叠加长存/长鑫国产化率从 15\%→25\% 的 S 曲线跃升，预计公司 2026E 收入 520--550 亿元，归母 68--71 亿元。平台化能力 + 服务体系 + 长存/长鑫深度绑定三大护城河，即使 HBM 竞争格局变化，设备 capex 确定性不变。2026E PE 41x。\textbf{核心风险}：研发投入持续高强度侵蚀短期利润；长存/长鑫扩产节奏受 ASML DUV 交付约束。

\textbf{中微公司 (688012.SH) — 刻蚀全球竞争力 + TSV HBM 入口}
\par 中微 CCP (电容耦合等离子体) 介质刻蚀在 3D NAND 领域全球竞争力已获长存验证 (长存介质刻蚀 约20\% 份额)。更重要的是 TSV (硅通孔) 刻蚀——HBM 的 12 层堆叠需要高精度 TSV 刻蚀，中微已进入 TSMC 先进封装供应链验证。2026E 营收 175--190 亿元，归母 30 亿元，2026E PE 48x。TSV 刻蚀的 HBM 入口价值尚未被市场充分定价。

\subsection{核心组合 Tier 2：高确定性 (7 只)}

\textbf{拓荆科技 (688072.SH)}：ALD/CVD 薄膜沉积双平台，长存 ALD 份额 >10\%。2026E PE 55x，增速 >45\%，PEG 合理。
\par \textbf{沪硅产业 (688126.SH)}：国内 12 寸硅片龙头 (全球前 6)，国内份额 约70\%。14nm 级硅片验证 2026 关键节点。
\par \textbf{深科技 (000021.SZ)}：Kingston 核心模组代工商，沛顿存储封测 DDR5 代工订单 2026E 翻倍。估值最低 (24x PE) 的确定性标的。
\par \textbf{长电科技 (600584.SH) + 通富微电 (002156.SZ)}：先进封装双龙头。长电 XDFOI 对标 CoWoS，通富 AMD 深度绑定。AI 先进封装收入实际贡献当前 <3\%，但 2027H2 后弹性巨大。
\par \textbf{江波龙 (301308.SZ)}：企业级 SSD + Lexar 品牌全球化。2026Q1 归母 38.62 亿 (YoY +2644\%) 反映存储涨价弹性。26E PE 仅 18x，估值最具吸引力。
\par \textbf{兆易创新 (603986.SH)}：NOR 全球前三 + 自研 DDR5 小批量量产。利基 DRAM 涨价弹性最大的 A 股标的。

\subsection{卫星/主题层 (弹性+事件驱动)}

\textbf{南大光电 (300346.SZ)}：ArF 光刻胶长存验证通过为核心催化 (2026H2)，一旦通过估值重估空间 >50\%。\par
\textbf{华天科技/北京君正/芯原股份}：AI 敞口实质较小或估值偏高，归入观察池。

\vspace{0.4em}
\noindent\textbf{投资含义}：标的分层的核心结论是 A 股存储链呈现「头部集中、长尾分散」特征——澜起 + 设备三强 + 模组双雄 6 只标的合计贡献链上约 65\% 的 AI 敞口市值。\textbf{配置策略应聚焦确定性最高的 6 只 (核心 Tier 1+2) 占比 75\%+}，剩余 25\% 分配给弹性卫星 (长电/通富/兆易) 和主题催化 (南大光电)。避免追逐低确定性长尾标的的估值波动。
